\subsection{原始6乗根候補の三次normによる非整域性}
\label{subsec:h3-order2-22-primitive6-nondomain}
\begin{theorem}\label{thm:h3-order2-22-primitive6-nondomain}
前節の$\lambda^2-\lambda+1=0$について、$AB\ne0$なら
\[
 R=B_{\lambda,\alpha,\beta}/(A x^3+B y^3)
\]
は全$\alpha,\beta$で整域ではない。
従って対角curved chart全体のH3部分は、$\lambda=-1,\alpha=\beta=0$の既知符号族だけである。
\end{theorem}
\begin{proof}
$R$が整域だと仮定して矛盾を導く。
$u^6,x^6$は中心であり、$u,x^3$は$R$内で非零である。
従ってそれらを逆転した環も整域となる。
まず$u$を逆転した次数零部分に
\[
 X=xu^{-1},\qquad Y=yu^{-1},\qquad V=vu^{-1},\qquad\omega=\lambda^2
\]
を取る。元の六関係から厳密に
\begin{gather*}
 VX=\lambda XV,\quad VY=\lambda^{-1}YV,\\
 XY-\omega YX=\lambda\alpha-\lambda\beta V^2,\qquad
 A X^3+B Y^3=0
\end{gather*}
を得る。
最後の式では$x^3u^{-3}=-X^3,y^3u^{-3}=-Y^3$を使った。

\[
 \delta=\frac{\lambda\alpha}{\omega-1}
       =\frac{1-2\lambda}{3}\alpha,\qquad
 \gamma=\frac{\lambda\beta}{\omega^{-1}-\omega}
       =\frac{\lambda-2}{3}\beta,
 \qquad C=YX+\delta+\gamma V^2
\]
と置くと
\[
 VC=CV,\qquad XC=\omega CX,\qquad YC=\omega^{-1}CY.
\]
切断前の$u$-chartは、可換基底$k[C,V]$の一般化Weyl代数である。
具体的には
\[
 \sigma(C)=\omega C,\quad\sigma(V)=\lambda^{-1}V,\qquad
 a_0=C-\delta-\gamma V^2,
\]
\[
 YX=a_0,\quad XY=\sigma(a_0),\quad
 Xf=\sigma(f)X,\quad Yf=\sigma^{-1}(f)Y
\]
となる。
$C,V$の代数的独立性とこの表示の単射性は、元のPBW基底における
$C^iV^j$の最高$x,y$次数と、$X$逆転後のLaurent標準語で確認できる。

$\sigma^6=1$、$\sigma^3(a_0)=a_0$であり、三次normは
\[
 P=\prod_{j=0}^2\sigma^j(a_0)
  =C^3-\delta^3-\gamma^3V^6-3\delta\gamma CV^2.
\]
これは非零かつ$\sigma$不変である。
$X$を逆転すれば$Y=a_0X^{-1}$、$Y^3=PX^{-3}$となる。
$r^3=-B/A$を選ぶと、切断関係は厳密に
\[
 X^6=r^3P
\]
へ移る。
このmonic関係のLaurent商は$k[C,V,P^{-1}]$上で
$1,X,\ldots,X^5$を基底に持つので、可換基底は切断後も単射である。
可換基底の非零元の$\sigma$-normは中心であり、それらを逆転して
$K=k(C,V)$へ移れる。整域の仮定の下ではこの局所化も整域でなければならない。
得られる環は
\[
 L=K[X,X^{-1};\sigma]/(X^6-r^3P)
\]
で、$K$上6次元の同じmonic基底を持つ。

ところが、$\sigma^2$による三次normも
\[
 a_0\sigma^2(a_0)\sigma^4(a_0)=P
\]
である。従って
\[
 Z=(r a_0)^{-1}X^2\in L
\]
と置くと$Z^3=1$であり、
\[
 (Z-1)(Z^2+Z+1)=0.
\]
両因子は非零である。
実際、第一因子はmonic基底の次数0と2、第二因子は次数0と2と4に非零係数を持つ。
よって$L$は整域ではなく、矛盾する。
この議論は$\delta$や$\gamma$の非零性を仮定せず、全$\alpha,\beta$を同時に覆う。
\end{proof}

検証器は
\[
 \texttt{verify\_cubic\_order2\_22\_primitive6\_nondomain.py}
\]
である。Laurent語からのdehomogenization、一般化Weyl関係、二つの三次norm、
6基底の乗法と二つの非零因子の零積を厳密再生する。
三次normが実際の非零因子の零積を与えることが排除の核心である。
